\subsubsection{Subproceso 3.B: Confirmaci\'on y Registro Oficial}

El Subproceso 3.B tiene como prop\'osito formalizar, consolidar y asentar con car\'acter legal y definitivo la matr\'icula del estudiante en el SIGA. Sobre esa base ejecuta el cierre del proceso de matr\'icula regular, emite los documentos oficiales de acreditaci\'on (Reporte y Boleta de Matr\'icula), distribuye los padrones de clase a los docentes y reporta la n\'omina consolidada a los organismos del Estado, SUNEDU y MINEDU.

Este subproceso constituye el punto de cierre de la cadena de valor de la matr\'icula regular: transforma una carga acad\'emica validada administrativamente (la Pre-Ficha Aprobada o Subsanada emitida por el Subproceso 3.A) en un registro oficial con validez jur\'idica ante la instituci\'on y ante el Estado. Su ejecuci\'on correcta garantiza la trazabilidad del acto de matr\'icula, habilita los tr\'amites que dependen de la Boleta de Matr\'icula (becas, seguros, expedici\'on del carn\'e universitario) y consolida la informaci\'on estad\'istica que la universidad debe reportar. Asimismo, su registro electr\'onico encuentra respaldo normativo en la Disposici\'on Transitoria \'Unica del Reglamento de Matr\'icula (RR-0570-2022), que faculta a la dependencia central de registro a simplificar y automatizar los tr\'amites de matr\'icula.
